\documentclass[11pt,a4paper]{article}

% ---------- packages ----------
\usepackage[utf8]{inputenc}
\usepackage[T1]{fontenc}
\usepackage{booktabs}
\usepackage{graphicx}
\usepackage{amsmath}
\usepackage{hyperref}
\usepackage[margin=1in]{geometry}
\usepackage{natbib}
\usepackage{xcolor}

\newcommand{\PH}[1]{\textcolor{red}{\textbf{[#1]}}}

% ---------- metadata ----------
\title{Semantic Novelty Measurement for Automated\\
Research Evaluation}

\author{
  Autonomous AI Researcher \\
  PicoClaw Research \\
  \texttt{research@picoclaw.dev}
}

\date{\today}

% ==========================================================================
\begin{document}
\maketitle

% ---------- abstract ----------
\begin{abstract}
Autonomous research systems need quantitative novelty signals to
distinguish genuinely original contributions from incremental extensions.
We evaluate four automated novelty metrics---embedding distance from
corpus centroid, atypical reference combinations, topic cluster distance,
and LLM-based novelty judgment---plus an optimally weighted combined
metric, on a corpus of 100 recent AI/ML papers annotated by expert proxy
raters.  Each metric is assessed by its ability to discriminate
expert-labeled novel papers from incremental ones (target AUC $\geq 0.70$)
and by its rank correlation with 2-year citation counts (target Spearman
$\rho \geq 0.30$).  \PH{Summary of main results: best AUC, best
Spearman, combined metric performance, hypothesis outcome.}
These results inform the SUNFIRE scoring system (E06) and all downstream
optimization loops that rely on automated novelty assessment.
\end{abstract}

% ==========================================================================
\section{Introduction}
\label{sec:introduction}

A central challenge for autonomous AI research systems is measuring the
novelty of generated outputs.  Without a reliable novelty signal, an
optimization loop risks converging on well-trodden ideas that score
highly on feasibility and clarity but contribute nothing new to the
literature.  Worse, if the novelty metric is noisy or biased, the system
may systematically favor a particular style of output---verbose but
vacuous, or contrarian but implausible---rather than genuinely novel
research directions.

Prior work on measuring scientific novelty falls into two broad
categories: \emph{bibliometric} approaches that analyze citation
patterns and reference combinations~\citep{uzzi2013atypical,
arts2020natural}, and \emph{semantic} approaches that compute distances
in learned representation spaces~\citep{beaty2021automating,
grootendorst2022bertopic}.  Bibliometric methods have the advantage of
grounding in real-world scholarly behavior but require citation data that
may not be available for newly generated ideas.  Semantic methods can be
applied to any text but require validation against expert judgments.

In this work (Experiment~E03), we evaluate four novelty metrics spanning
both categories, plus a learned combination, on a corpus of 100 recent
AI/ML papers.  We contribute:

\begin{enumerate}
  \item A curated dataset of 100 AI/ML papers with expert proxy novelty
        annotations (binary novel/incremental labels and 1--7 Likert
        scores).
  \item Four individual novelty metrics: embedding distance from corpus
        centroid, atypical reference combinations, topic cluster distance,
        and LLM-based novelty judgment.
  \item An optimally weighted combined metric trained on the annotation
        labels.
  \item A systematic evaluation of all metrics against expert labels
        (AUC-ROC) and 2-year citation counts (Spearman correlation),
        with bootstrap confidence intervals.
\end{enumerate}

% ==========================================================================
\section{Related Work}
\label{sec:related}

\paragraph{Atypical combinations and novelty.}
\citet{uzzi2013atypical} analyzed 17.9 million papers in the Web of
Science and showed that the highest-impact papers combine highly
conventional pairings of referenced journals with unusual pairings---a
signature of atypical knowledge recombination.  This finding motivates
our atypical reference combinations metric, which measures the
statistical unusualness of a paper's reference pairings relative to a
background corpus.

\paragraph{Natural language novelty indicators.}
\citet{arts2020natural} developed text-based novelty measures using
natural language processing to identify novel scientific contributions.
They demonstrated that linguistic markers of novelty in abstracts
correlate with citation impact and patent citations, suggesting that
semantic features of text capture meaningful signals about research
originality.

\paragraph{Automated creativity assessment.}
\citet{beaty2021automating} used semantic distance methods to
automatically score divergent thinking tasks, showing that embedding
distance from a reference distribution correlates with expert creativity
ratings ($r \approx 0.80$).  We adapt this approach to measure the
semantic distance of a paper's abstract from the corpus centroid as a
proxy for topical novelty.

\paragraph{Topic modeling.}
\citet{grootendorst2022bertopic} introduced BERTopic, a neural topic
modeling approach that clusters document embeddings and extracts topic
representations.  We use topic modeling to assign papers to clusters and
measure each paper's distance from its assigned cluster centroid as a
within-topic novelty signal.

% ==========================================================================
\section{Methodology}
\label{sec:methodology}

\subsection{Paper Sampling}
\label{sec:sampling}

We sample 100 papers from Semantic Scholar using the following criteria:

\begin{itemize}
  \item Published in the AI/ML domain within the past 3 years.
  \item Have at least 2 years of citation data available.
  \item Represent a range of venues (top conferences, workshops, and
        preprints) to ensure variance in novelty.
  \item Include both highly cited and low-citation papers.
\end{itemize}

For each paper, we retrieve the title, abstract, reference list,
venue, publication date, and 2-year citation count from the Semantic
Scholar API.

\subsection{Expert Proxy Annotation}
\label{sec:annotation}

Three expert proxy raters are instantiated using Claude Opus 4 with
distinct personas:

\begin{enumerate}
  \item \textbf{Domain Expert}: Senior ML researcher evaluating
        theoretical and methodological novelty.
  \item \textbf{Applied Researcher}: Industry researcher evaluating
        practical novelty and differentiation from existing tools.
  \item \textbf{Interdisciplinary Reviewer}: Researcher at the
        intersection of AI and another field, evaluating cross-domain
        novelty.
\end{enumerate}

Each rater annotates all 100 papers with:
\begin{itemize}
  \item A binary label: \textbf{novel} or \textbf{incremental}.
  \item A 1--7 Likert novelty score (1 = entirely derivative,
        7 = highly original).
  \item A short rationale for the rating.
\end{itemize}

Inter-rater agreement is measured via Fleiss' kappa on the binary
labels.

\subsection{Novelty Metrics}
\label{sec:metrics}

\paragraph{Metric 1: Embedding distance from centroid.}
We compute TF-IDF vectors over paper abstracts, then measure the cosine
distance of each paper's vector from the corpus centroid.  Papers farther
from the centroid are scored as more novel.

\paragraph{Metric 2: Atypical reference combinations.}
Following~\citet{uzzi2013atypical}, we compute pairwise co-occurrence
frequencies of referenced venues across the corpus.  For each paper, we
calculate the median $z$-score of its reference-pair frequencies.  Lower
$z$-scores (more unusual combinations) indicate higher novelty.

\paragraph{Metric 3: Topic cluster distance.}
We cluster papers into $k$ topics using $k$-means on TF-IDF embeddings
and measure each paper's cosine distance from its assigned cluster
centroid.  Papers that sit at the boundary between clusters or far from
their cluster center receive higher novelty scores.

\paragraph{Metric 4: LLM novelty judgment.}
Two LLM judges (Claude Opus 4.6 and Claude Sonnet 4.6) independently
rate each paper's novelty on a 1--7 scale given the title, abstract, and
a brief description of the paper's main contribution.  Scores are
averaged across judges.

\paragraph{Combined metric.}
We compute a weighted combination of the four individual metrics, with
weights optimized via logistic regression on the binary expert labels
using leave-one-out cross-validation to avoid overfitting.

\subsection{Statistical Tests}
\label{sec:statistics}

\begin{itemize}
  \item \textbf{AUC-ROC} for each metric's ability to discriminate
        expert-labeled novel papers from incremental ones, with
        bootstrap 95\% CIs ($B = 10{,}000$).
  \item \textbf{Spearman's $\rho$} between each metric and 2-year
        citation counts, with permutation $p$-values
        ($10{,}000$ permutations).
  \item \textbf{Spearman's $\rho$} between each metric and mean
        human Likert scores.
  \item \textbf{Precision@$k$} ($k = 10, 20$) for identifying top
        expert-rated novel papers.
  \item \textbf{Fleiss' kappa} for inter-rater agreement on binary
        labels.
  \item \textbf{DeLong test} for comparing AUCs between metrics.
\end{itemize}

\subsection{Hypotheses}
\label{sec:hypotheses}

\begin{description}
  \item[H1:] At least one metric achieves AUC $\geq 0.70$ for
    discriminating expert-labeled novel vs.\ incremental papers.
  \item[H2:] At least one metric achieves Spearman $\rho \geq 0.30$
    with 2-year citation counts.
\end{description}

Both hypotheses are pre-registered before data analysis.

% ==========================================================================
\section{Results}
\label{sec:results}

\subsection{Inter-Rater Agreement}

\PH{Fleiss' kappa value and interpretation.  Agreement level
(slight/fair/moderate/substantial) per Landis \& Koch.}

\subsection{Individual Metric Performance}

Table~\ref{tab:metrics} reports the discrimination and correlation
performance of each novelty metric.

\begin{table}[ht]
\centering
\caption{Novelty metric performance.  AUC for binary discrimination,
Spearman $\rho$ for correlation with citation counts and human Likert
scores.  Bootstrap 95\% CIs in brackets.}
\label{tab:metrics}
\begin{tabular}{lccccc}
\toprule
\textbf{Metric} & \textbf{AUC} & \textbf{95\% CI} &
  \textbf{$\rho$ (citation)} & \textbf{$\rho$ (human)} &
  \textbf{P@10} \\
\midrule
Embedding distance   & \PH{auc} & \PH{ci} & \PH{rho} & \PH{rho} & \PH{p10} \\
Atypical references  & \PH{auc} & \PH{ci} & \PH{rho} & \PH{rho} & \PH{p10} \\
Topic distance       & \PH{auc} & \PH{ci} & \PH{rho} & \PH{rho} & \PH{p10} \\
LLM judgment         & \PH{auc} & \PH{ci} & \PH{rho} & \PH{rho} & \PH{p10} \\
Combined             & \PH{auc} & \PH{ci} & \PH{rho} & \PH{rho} & \PH{p10} \\
\bottomrule
\end{tabular}
\end{table}

\PH{Narrative description of which metrics performed best and worst,
notable patterns in the data.}

\subsection{ROC Curves}

Figure~\ref{fig:roc} shows the ROC curves for all five metrics.

\begin{figure}[ht]
\centering
\PH{ROC curve figure: includegraphics roc\_curves.pdf}
\caption{ROC curves for all novelty metrics.  The dashed diagonal
represents random classification.}
\label{fig:roc}
\end{figure}

\subsection{Metric Correlations}

\PH{Inter-metric correlation matrix.  Discussion of which metrics
capture overlapping vs.\ complementary signals.}

\subsection{Citation Correlation}

\PH{Scatter plots or discussion of the relationship between novelty
scores and 2-year citation counts.  Note any non-linearity.}

% ==========================================================================
\section{Discussion}
\label{sec:discussion}

\subsection{Hypothesis Evaluation}

\PH{H1 evaluation: which metrics achieved AUC >= 0.70.  Pass/fail
with exact values and CIs.}

\PH{H2 evaluation: which metrics achieved Spearman rho >= 0.30 with
citations.  Pass/fail with exact values and p-values.}

\PH{Overall assessment: hypotheses supported or not supported.}

\subsection{Key Findings}

\PH{Discussion of which metric category (bibliometric vs.\ semantic
vs.\ LLM) performs best and why.  Value of the combined metric over
individual metrics.  Practical recommendations for which metric to
use in the autonomous research pipeline.}

\subsection{Limitations}

\begin{itemize}
  \item \textbf{Expert proxy rather than human experts.}  Ground truth
    annotations come from LLM-based expert proxies rather than genuine
    domain experts.  This introduces potential shared biases between
    annotator and metric models, particularly for the LLM judgment
    metric.
  \item \textbf{Single domain.}  All papers are from AI/ML; novelty
    metrics may behave differently in other scientific domains where
    the structure of knowledge is different.
  \item \textbf{Citation as novelty proxy.}  Citation counts reflect
    impact rather than novelty directly.  A paper can be highly cited
    for reasons unrelated to novelty (e.g., providing a useful
    benchmark or survey).
  \item \textbf{TF-IDF limitations.}  TF-IDF embeddings capture lexical
    but not deep semantic similarity.  Transformer-based embeddings
    (e.g., Specter) may improve embedding-based metrics.
  \item \textbf{Sample size.}  100 papers provides moderate statistical
    power; effects smaller than $d \approx 0.3$ may not be detected.
  \item \textbf{Temporal confound.}  Older papers in the sample have had
    more time to accumulate citations, potentially confounding the
    citation correlation analysis.
\end{itemize}

\subsection{Implications for Downstream Experiments}

\begin{itemize}
  \item \textbf{E06 (SUNFIRE scoring):}  The best-performing novelty
    metric (or the combined metric) can be integrated as the novelty
    component of the SUNFIRE scoring function.
  \item \textbf{E04 (Gap detection):}  Novelty scores can be used to
    weight research gaps by the novelty of potential contributions.
  \item \textbf{Optimization loops (E09--E15):}  A validated novelty
    metric enables fitness functions that reward genuine originality
    rather than superficial variation.
\end{itemize}

% ==========================================================================
\section{Conclusion}
\label{sec:conclusion}

\PH{Summary of results: best metric, combined metric performance,
hypothesis outcomes.  Practical recommendation for the autonomous
research pipeline.  Brief statement of what this enables for
downstream experiments.}

% ==========================================================================
\bibliographystyle{plainnat}

\begin{thebibliography}{6}

\bibitem[Uzzi et~al.(2013)]{uzzi2013atypical}
Brian Uzzi, Satyam Mukherjee, Michael Stringer, and Ben Jones.
\newblock Atypical combinations and scientific impact.
\newblock \emph{Science}, 342(6157):468--472, 2013.

\bibitem[Arts and Veugelers(2020)]{arts2020natural}
Sam Arts and Reinhilde Veugelers.
\newblock Natural language processing to identify the creation and impact of
  new technologies in patent text: Code, data, and new measures.
\newblock \emph{Research Policy}, 49(10):104144, 2020.

\bibitem[Beaty and Johnson(2021)]{beaty2021automating}
Roger~E. Beaty and Dan~R. Johnson.
\newblock Automating creativity assessment with {SemDis}: An open platform for
  computing semantic distance.
\newblock \emph{Behavior Research Methods}, 53(2):757--780, 2021.

\bibitem[Grootendorst(2022)]{grootendorst2022bertopic}
Maarten Grootendorst.
\newblock {BERTopic}: Neural topic modeling with a class-based {TF-IDF}
  procedure.
\newblock \emph{arXiv preprint arXiv:2203.05794}, 2022.

\bibitem[Lu et~al.(2024)]{lu2024aiscientist}
Chris Lu, Cong Lu, Robert~Tjarko Lange, Jakob Foerster, Jeff Clune, and
David Ha.
\newblock The {AI} {S}cientist: Towards fully automated open-ended scientific
  discovery.
\newblock \emph{arXiv preprint arXiv:2408.06292}, 2024.

\bibitem[Efron and Tibshirani(1993)]{efron1993bootstrap}
Bradley Efron and Robert~J. Tibshirani.
\newblock \emph{An Introduction to the Bootstrap}.
\newblock Chapman \& Hall, 1993.

\end{thebibliography}

\end{document}
